\documentclass[10pt,a4paper]{amsart}
\usepackage[margin=19mm]{geometry}
\usepackage{amsmath,amssymb,mathtools}
\usepackage[scr=rsfs,cal=euler]{mathalfa}
\usepackage{xfrac}
\usepackage[unicode,hidelinks,bookmarks=false]{hyperref}
\newcommand{\ie}{i.e.\ }
\newcommand{\cf}{cf.\ }
\newcommand{\resp}{respectively\ }
\newcommand{\Subsection}{\subsection}
\providecommand{\nobreakdash}{\nobreak\hbox{-}}
\setlength{\emergencystretch}{2em}
\multlinegap0pt
\usepackage{amssymb}
\usepackage{mathtools}
\usepackage{bm}
\usepackage{float}
\usepackage{xcolor}
\newtheorem{theorem}{Theorem}
\title{Identification and Inference in Nonlinear Dynamic Network Models}
\author{Diego Vallarino}
\date{Source: arXiv:2604.04961v1 (CC BY 4.0)}
\begin{document}
\maketitle

\noindent\emph{Source and licence.} Identification and Inference in Nonlinear Dynamic Network Models. Version arXiv:2604.04961v1 (CC BY 4.0), distributed by arXiv under the Creative Commons Attribution 4.0 International licence, http://creativecommons.org/licenses/by/4.0/, as recorded in the arXiv licence metadata for this exact version and shown on the abstract page https://arxiv.org/abs/2604.04961v1. Full licence notice: \texttt{licenses/arxiv-2604.04961-CC-BY-4.0.txt}.\par\smallskip

\noindent\textit{Editorial scope.} Counted unit: the article's theorem eq:4.exchangeable (source0.tex lines 639--661). The article's complete original proof of it is delivered with it (source0.tex lines 663--677, one \texttt{proof} environment of 15 lines). No mathematics is written, repaired or re-derived: the statement and the proof are the article's own lines, in the article's own notation, and no retained line is substituted or re-flowed.  No further source-local context is needed: every cross-reference of every retained line resolves inside the two delivered spans.  Nothing was omitted from any retained span, so the package carries no omission record.  No retained line contains a citation, so the wrapper carries no bibliography and no bibliography entry.  No retained line contains a figure environment, an image-inclusion command or drawing code, so no image is delivered and the package declares no asset dependency.  Editorial formatting only, in addition: an AMS \texttt{amsart} preamble that restates the article's own declarations and macros used by the retained lines, namely the environments \texttt{theorem} on separate counters, together with the standard packages those lines need.  The retained environments print in this document as Theorem 1 in order of appearance, whereas the article prints the same items with the numbering of its own full document.  The complete native source of the article as distributed in the arXiv e-print for this exact version is bundled unchanged as \texttt{sources/197922/source0.tex}.
\begin{theorem}[Non-identification under exchangeability]
Suppose

\begin{equation}
\Sigma
=
\sigma^2 I
+
\tau^2 \mathbf{1}\mathbf{1}',
\label{eq:4.exchangeable}
\end{equation}

and

\begin{equation}
\Omega
=
\omega^2 I.
\label{eq:4.exchangeable2}
\end{equation}

Then the interaction matrix $A$ is not identified.
\end{theorem}

\begin{proof}
Under exchangeability, $\Sigma$ commutes with any permutation matrix.
Therefore, any interaction matrix of the form

\begin{equation}
A'
=
P A P'
\label{eq:4.permutation}
\end{equation}

with $P$ permutation, generates the same covariance matrix.
Since the observable moments depend only on $\Sigma$, the interaction
matrix cannot be uniquely recovered.
\end{proof}

\end{document}
